%%%%%%%%%%%%%%%%%%%%%%%%%%%%%%%%%%%%%%%%%%%%%%%%%%%%%%%%%%%%%%%%%%%%%%%%%%%%%%%
% SECTION 4: History
%%%%%%%%%%%%%%%%%%%%%%%%%%%%%%%%%%%%%%%%%%%%%%%%%%%%%%%%%%%%%%%%%%%%%%%%%%%%%%%
\section{History}

% The AI Winters
\begin{frame}{The AI Winters}
    \begin{table}
        \centering
        \small
        \begin{tabular}{@{}ll@{}}
            \toprule
            \textbf{Period} & \textbf{Event} \\ \midrule
            1950s--60s & Symbolic AI, first wave of optimism \\
            1970s & \textbf{AI Winter 1} — symbolic approaches hit limits \\
            1980s & Expert systems — encode human knowledge explicitly \\
            1990s & \textbf{AI Winter 2} — expert systems too brittle \\
            1986 & Hinton's backprop paper — \textit{largely ignored} \\
            1990s & LeCun's CNNs deployed — \textit{largely ignored} \\
            2000s & SVMs and kernel methods dominate mainstream \\ \bottomrule
        \end{tabular}
    \end{table}
    
    \vspace{1em}
    Deep learning survived two cycles of hype and abandonment
    
    \textbf{The authors lived through both winters, continued regardless}
\end{frame}

% The 2006 Breakthrough - Problem
\begin{frame}{The 2006 Breakthrough: The Problem}
    \textbf{Training deep networks was extremely difficult:}
    
    \vspace{1em}
    \begin{itemize}
        \item No ReLU — stuck with sigmoid/tanh
        \item Poor weight initialization
        \item No dropout regularization
        \item Gradients vanished before reaching early layers
        \item Limited labeled data
        \item Limited compute power
    \end{itemize}
\end{frame}

% Boltzmann Machines
\begin{frame}{Boltzmann Machines}
    % TODO: Add Boltzmann Machine structure diagram
    % - Visible units (bottom)
    % - Hidden units (top)
    % - Connections between layers
    % - Note: connections are bidirectional/symmetric
    
    \begin{center}
        \textit{[Diagram: Boltzmann Machine structure]}
    \end{center}
    
    \vspace{1em}
    \textbf{Hinton \& Sejnowski, 1985}
    \begin{itemize}
        \item Stochastic recurrent neural network
        \item Can learn internal representations without labels
        \item \textbf{Restricted Boltzmann Machines (RBMs):} no intra-layer connections
    \end{itemize}
    
    \vspace{0.5em}
    Foundation for the 2006 breakthrough
\end{frame}

% The 2006 Breakthrough - Solution
\begin{frame}{The 2006 Breakthrough: The Solution}
    \textbf{Greedy layer-wise pretraining:}
    
    \vspace{1em}
    \begin{enumerate}
        \item Train each layer independently using \textbf{unlabeled data}
        \item Use Restricted Boltzmann Machines
        \item Stack the pretrained layers
        \item Fine-tune end-to-end with labeled data
    \end{enumerate}
    
    \vspace{1em}
    \textbf{Why it mattered:}
    \begin{itemize}
        \item Proved deep nets could learn features \textbf{without labels}
        \item Proved pursuing depth was worthwhile at all
    \end{itemize}
    
    \vspace{1em}
    \begin{quote}
        \small\textit{Note: the paper attributes this to ``a group at CIFAR.'' Hinton and LeCun were that group.}
    \end{quote}
\end{frame}

% First Commercial Payoff
\begin{frame}{First Commercial Payoff: Speech Recognition}
    \begin{itemize}
        \item Adopted by major industry groups \textasciitilde2009
        \item Deployed in Android phones by 2012
        \item Compounded by rapidly improving GPU hardware
    \end{itemize}
    
    \vspace{1em}
    \textbf{GPUs:} Originally for graphics/gaming
    \begin{itemize}
        \item Massively parallel architecture
        \item Perfect for matrix operations in neural networks
        \item Training times dropped from weeks to hours
    \end{itemize}
    
    \vspace{1em}
    Once training innovations arrived, pretraining became unnecessary for large labeled datasets
\end{frame}
